\section{D. Profinite extensions and the limits of the realization problem}

We write $P=\widehat\pi_0$, where $\pi_0=\pi_1(U,p)$ for a connected orientable surface of finite type with $\chi(U)<0$. All maps between profinite groups in this section are continuous. A lift of an outer action means a homomorphism inducing the prescribed outer homomorphism, with its labels in the operator group retained. For a finite group $G$ acting faithfully on $U$, $E$ denotes the group of lifts to the universal cover, not a mapping class group. The completion map $E\to\widehat E$ is injective: $E$ is a finitely generated linear hyperbolic orbifold group and is residually finite. We use $E^+$ for its orientation-preserving subgroup. The passages answered here are \S31, folios 198--199, scans 391--394; the inserted $G=G^+$ on scan 391 is material.

\begin{claim}{D.1}{Proved}{Exact completion of the lift extension}
The completion of $1\to\pi_0\to E\xrightarrow qG\to1$ is the exact sequence
\[
1\longrightarrow P\longrightarrow\widehat E\xrightarrow{\widehat q}G\longrightarrow1.
\]
It is the extension attached to the completed outer action. An invariant discretization therefore supplies a canonical comparison of extensions, with identity maps on $P$ and $G$.
\end{claim}
\begin{proof}
Every subgroup of finite index in $\pi_0$ contains a subgroup normal and of finite index in $E$: take its core in $E$, using $[E:\pi_0]<\infty$. Thus $\pi_0$ has its full profinite topology induced from $E$, and its closure is $P$. The kernel of $\widehat q$ is that closure, since $\pi_0=q^{-1}(1)$ is open. Conjugation in the completed extension is the continuous extension of conjugation in the dense subgroup $E$. The uniqueness of a center-free kernel extension, proved in Part A, identifies this extension with the pullback of $\Aut(P)\to\Out(P)$. The identification fixes the named kernel and quotient, so the assertion of canonicity has no unmentioned conjugator.
\end{proof}

\begin{claim}{D.2}{Proved}{Finite subgroups of completed surface orbifold groups, including reflections}
Every finite subgroup of $\widehat E$ is conjugate to a finite subgroup of $E$. Two finite subgroups of $E$ conjugate in $\widehat E$ are conjugate in $E$. These finite subgroups are cyclic rotation groups or dihedral groups, with a single reflection included as the dihedral case of order two.
\end{claim}
\begin{proof}
We specify the two inputs and the reduction. Bridson--Conder--Reid, Theorem 5.1 \cite{R-BCR}, applies to the finitely generated Fuchsian group $E^+$: every finite subgroup of $\widehat{E^+}$ is conjugate into $E^+$. Boggi--Zalesskii, Theorem 6.3 \cite{R-BZ}, gives conjugacy and normalizer comparison for finite $p$-subgroups of $E$:
\[
N_{\widehat E}(A)=\overline{N_E(A)}.
\]
Here its hypotheses hold because $E$ has a finite-index surface group (a free group in the noncompact case), hence is virtually compact special; its good-group and finiteness hypotheses are precisely those in Theorem 6.3. We use only this prime-power theorem, not an assertion obtained by commuting arbitrary intersections with closure.

Put $K^+=K\cap\widehat{E^+}$ for a finite $K\leq\widehat E$. The first input makes $K^+$ cyclic. When $K^+=1$, the orientation character embeds $K$ in $C_2$, and the prime-power input handles $K$. Otherwise choose the unique subgroup $A$ of order $p$ in $K^+$, for one prime $p\mid |K^+|$. It is characteristic in $K^+$ and normal in $K$. Conjugate $A$ into $E^+$. A nontrivial elliptic rotation of the hyperbolic plane has one fixed point $x$. Its normalizer in the discrete group $E$ stabilizes $x$, and conversely every element of $E_x$ normalizes the unique subgroup of order $p$ in its cyclic rotation subgroup. Thus $N_E(A)=E_x$ is finite cyclic or dihedral. The displayed comparison gives $N_{\widehat E}(A)=E_x$, so $K\leq E_x\leq E$.

For injectivity, suppose $K,L\leq E$ and $tKt^{-1}=L$, with $t\in\widehat E$. The orientation character gives $tK^+t^{-1}=L^+$. In the case $K^+=1$ apply prime-power conjugacy separation. Otherwise choose $A\leq K^+$ as above and $B=tAt^{-1}\leq L^+$. Prime-power conjugacy separation supplies $e\in E$ with $eAe^{-1}=B$. Now $e^{-1}t$ normalizes $A$, hence belongs to $N_E(A)\subset E$. Therefore $t\in E$, after this adjustment. This proves the required conjugacy assertion. Finally finite discrete hyperbolic-plane isometry groups stabilize a point and embed in $O(2)$; the listed cyclic and dihedral possibilities follow.
\end{proof}

\begin{claim}{D.3}{Proved}{The comparison of full and partial section classes}
For every subgroup $H\leq G$, completion induces a bijection
\[
\operatorname{Sec}(E_H\to H)/\pi_0
\ \longrightarrow\
\operatorname{Sec}(\widehat E_H\to H)/P,
\qquad E_H=q^{-1}(H).
\]
Both sets are finite. In particular the first conjecture on scan 393 holds for geometrically realized actions, with orientation-reversing operators allowed.
\end{claim}
\begin{proof}
We prove the change from total-group conjugacy to kernel conjugacy, since omission of this step changes the question. A profinite section has a finite image $K$ mapping isomorphically to $H$. Apply D.2 to $E_H$, and choose $t\in\widehat E_H$ with $tKt^{-1}\leq E_H$. Choose $e\in E_H$ with $q(e)=\widehat q(t)$. Then $e^{-1}t\in P$, and conjugation by this element puts $K$ in $E_H$. The restriction of $q$ identifies the resulting subgroup with $H$, so it defines a discrete section with the correct labels.

Conversely, let discrete section images $K,L$ be conjugate by an element of $P$. D.2, applied in $E_H$, gives an $e\in E_H$ with $eKe^{-1}=L$. Let $l\in L$ be the unique element satisfying $q(l)=q(e)$. Then $l^{-1}e\in\pi_0$ still conjugates $K$ to $L$. Because the projection is injective on each section image, this equality of subgroups is equality of the labelled section homomorphisms after conjugation. This proves injectivity.

There are finitely many $E_H$-conjugacy classes of finite subgroups: they occur in the finite list of cone stabilizers, corner stabilizers, and mirror components of a finite-type quotient orbifold, and each finite stabilizer has finitely many subgroups. Each such class splits into at most $[E_H:\pi_0]=|H|$ kernel-conjugacy classes. The asserted finiteness follows.
\end{proof}

\begin{claim}{D.4}{Proved}{The exact profinite fixed subgroup of a section}
Let $s:H\to E_H$ be a nontrivial discrete section. Its fixed subgroup $P^{s(H)}=C_P(s(H))$ is trivial when $s(H)$ contains a nontrivial rotation. For a reflection $s(H)=\langle t\rangle$, it is $1$ for a noncompact real component and is the closure of the primitive oval subgroup, isomorphic to $\widehat\Z$, for a compact real component. The same assertions hold for every profinite section by D.3.
\end{claim}
\begin{proof}
For a rotation choose $A$ of prime order inside $s(H)^+$. D.2's proof gives $N_{\widehat E}(A)=E_x$, a finite group. Thus $C_P(s(H))\leq E_x\cap P=1$, because $P$ is torsion-free. For a reflection the prime-power centralizer comparison gives
\[
C_{\widehat E}(t)=\overline{C_E(t)},\qquad
C_P(t)=\overline{C_{\pi_0}(t)}.
\]
The second equality uses that $P$ is open and $P\cap E=\pi_0$: the intersection of the closure of a subgroup with an open subgroup is the closure of their intersection. On the universal cover, $C_{\pi_0}(t)$ is the group of deck transformations preserving and commuting with reflection in its fixed line. It acts freely on that line. Its quotient is the corresponding real component of $U$. The deck group is trivial for a line and infinite cyclic for a circle; the circle is an essential primitive simple curve. Its induced profinite topology is full, so its closure is $\widehat\Z$.

This proves, rather than assumes, the correction to the expected answer in Part D(a): a dihedral section with a nontrivial rotation has fixed subgroup $1$. The value $\widehat\Z$ belongs to a reflection in an oval. It is not the fixed subgroup of a larger dihedral stabilizer containing that reflection.
\end{proof}

\begin{claim}{D.5}{Refuted}{Unique maximal sections require a rotation}
A finite subgroup of $\widehat E$ containing a nontrivial rotation is contained in a unique maximal finite subgroup. The assertion for every nontrivial finite subgroup is false, already for type $(0,3)$ with reflections.
\end{claim}
\begin{proof}
For the positive assertion conjugate the subgroup into $E$ using D.2. Its rotation part has a subgroup $A$ of prime order. Every finite overgroup has cyclic orientation-preserving part, so normalizes this same $A$. All finite overgroups therefore lie in the finite group $N_{\widehat E}(A)=E_x$, which is itself the unique maximal one.

For the counterexample use the $(2,3,\infty)$ reflection group
\[
W=\langle s_0,s_1,s_2\mid s_i^2=1,
(s_0s_1)^2=(s_1s_2)^3=1\rangle
       = V_4*_{C_2}S_3.
\]
Map it onto $S_3\times C_2$, with $r^3=a^2=\tau^2=1$, $ara=r^{-1}$ and $\tau$ central, by
$s_0\mapsto\tau$, $s_1\mapsto a\tau$, $s_2\mapsto ra\tau$.
The vertex groups inject. The kernel is an orientable torsion-free group of index $12$, Euler characteristic $-1$, with three cusps; hence it is $\pi_{0,3}$. The two maximal finite subgroups
$\langle s_0,s_1\rangle$ and $\langle s_1,s_2\rangle$ contain the same reflection $s_1$ and have distinct images in the quotient. They remain distinct and are not kernel-conjugate after completion. Their generated subgroup is $W$, so no finite overgroup contains both. Part B supplies the finite certificate and the geometric kernel computation.
\end{proof}

\begin{claim}{D.6}{Proved}{A well-defined transfer of the geometric rotation invariant}
For a realized $E$, its rotation invariant on finite-order elements transfers uniquely to the torsion of $\widehat E$ by conjugating each torsion element into $E$. This assertion uses the orientation module, so that conjugation by an orientation-reversing element acts by $-1$ on its values. Its values embed canonically in $\widehat T\otimes_{\Z}\Q/\Z$, where $\widehat T=T\otimes_{\Z}\widehat\Z$.
\end{claim}
\begin{proof}
Use D.2 for existence. Explicitly, for torsion $h\in\widehat E$ choose $t\in\widehat E$ with $u=tht^{-1}\in E$ and set
\[
 \mu_{\widehat E}(h)=\varepsilon(t)\mu_E(u),
\]
where $\varepsilon:\widehat E\to\{\pm1\}$ is the orientation character, acting on $\widehat T$. Suppose $u,v\in E$ and $v=tut^{-1}$ in $\widehat E$. For a rotation apply the proof of D.2 to an order-$p$ subgroup of $\langle u\rangle$: first align it by discrete conjugation, then put the remaining conjugator in the finite discrete stabilizer. Thus the elements themselves, and not just their cyclic subgroups, are discretely conjugate, with the same orientation character of the conjugator. For reflections apply prime-power element conjugacy separation, or D.2 because an order-two subgroup has a unique nonidentity element. This establishes independence with the indicated action on $T$. A rotation of order $n$ has a value in $T/nT\subset T\otimes\Q/\Z$; the identification $T/nT=\widehat T/n\widehat T$ defines the claimed inclusion.
\end{proof}

% Replacement for D.7.  The bibliography entry bWilkes2017Seifert is supplied
% in d_mu_references.bib.  This fragment has no global preamble.
\begin{claim}{D.7}{Proved}{Equivariance of the profinite rotation invariant}
Let $E$ be the lift group of a finite action on an orientable hyperbolic
surface of finite type, and put $P=\widehat\pi\triangleleft\widehat E$.
Every automorphism $\alpha$ of $\widehat E$ preserving $P$ and its indexed
peripheral structure satisfies
\[
 \mu_{\widehat E}(\alpha(u))=\alpha_T\bigl(\mu_{\widehat E}(u)\bigr)
 \qquad (u\in\widehat E\text{ of finite order}).
\]
Here $\alpha_T$ is the induced automorphism of the orientation line
$\widehat T$ defined in A.2, and $\mu_{\widehat E}$ is the invariant of
D.6.  The assertion holds for compact and noncompact surfaces, and for
actions containing orientation-reversing elements.
\end{claim}
\begin{proof}
Choose an orientation, so that $\alpha_T$ is multiplication by some
$\lambda\in\widehat\Z^\times$.  Conjugation by an element of
$\widehat E$ acts on $\widehat T$ by its geometric orientation sign.
Naturality of this action and commutativity of the automorphism group of a
rank-one module show that $\alpha$ preserves that sign.  Thus $\alpha$
restricts to an automorphism of $\widehat B$, where $B=E^+$ is the
orientable Fuchsian subgroup.  A nonidentity negative torsion element is a
reflection and has invariant $1/2$; multiplication by a profinite unit
fixes this value.  It remains to treat torsion in $\widehat B$.

We first make explicit the behavior of orientation under finite covers.
For a degree-$d$ connected finite cover of oriented surfaces, the map on
orientation homology takes the fundamental class upstairs to $d$ times
the fundamental class downstairs.  For punctured surfaces this follows
directly in A.2's kernel model: the peripheral factors above any one
puncture map by their positive local degrees, whose sum is $d$.
For closed surfaces it is the degree formula for fundamental classes.
The same formulas hold over $\widehat\Z$, either by this peripheral
matrix or by the natural finite-coefficient cohomology comparison.
Consequently, if $\alpha$ carries an open surface subgroup $V$ to $V'$,
the multipliers on their orientation lines and on $P$ satisfy
$d\lambda_V=d\lambda$.  Multiplication by a positive integer is
injective on $\widehat\Z=\prod_p\Z_p$, so $\lambda_V=\lambda$.
This proves the required compatibility with further finite covers; it
does not divide by $d$ in a finite coefficient ring.

\emph{The noncompact case.}
Let $L$ be a peripheral subgroup of $\pi$, and let $K$ be the primitive
cusp subgroup of $B$ containing it.  Write $K=\langle w\rangle$ and
$L=\langle w^n\rangle$, with compatible positive generators.
We have
\[
 N_{\widehat B}(\overline L)=\overline K\simeq\widehat\Z.
\]
Here is the reduction from Part A, rather than an assumption that arbitrary
normalizers commute with completion.  For an element of the normalizer,
choose a representative in $B$ of its coset modulo $P$.  That coset
stabilizes the profinite peripheral class.  Distinct discrete peripheral
classes remain distinct in $P$, so the same coset stabilizes the discrete
peripheral class.  Adjust its representative by an element of $\pi$ to
put it in $N_B(L)=K$.  The remaining factor belongs to
$N_P(\overline L)=\overline L$, as proved in Part A.  This proves the
displayed equality.

The loop condition on $\alpha|_P$ sends a positive generator of
$\overline L$ to a $P$-conjugate of the $\lambda$th power of the
corresponding generator of another peripheral subgroup $\overline L'$.
After making this conjugation, $\alpha$ sends $\overline K$ to
$\overline K'$ and has the form $w\mapsto(w')^a$.
The indices
$[\overline K:\overline L]=[\overline K':\overline L']=n$
are equal because $\alpha$ preserves $P$.  Thus $na=n\lambda$ in
$\widehat\Z$, and hence $a=\lambda$.  Every primitive orbifold cusp
therefore has precisely the same multiplier $\lambda$.

Write the signature of $B$ as $(h;m_1,\ldots,m_r;s)$ with $s\geq1$.
Choose cone generators $c_j$, of orders $m_j$, and cusp generators $w_i$
in a common orientation.  The abelianized presentation is
\[
 \widehat B^{\mathrm{ab}}
 =\left(\widehat\Z^{2h}\oplus
       \bigoplus_{j=1}^r\Z/m_j\,c_j\oplus
       \bigoplus_{i=1}^s\widehat\Z\,w_i\right)
       \big/\left\langle\sum_jc_j+\sum_iw_i\right\rangle.
\]
D.2 applied to $B$ makes $\alpha$ permute the maximal cone subgroups.
Hence, up to conjugacy in $\widehat B$,
$\alpha(c_j)=c_{\tau(j)}^{k_j}$, where
$m_{\tau(j)}=m_j$ and $k_j$ is a unit modulo $m_j$.
The preceding paragraph gives
$\alpha(w_i)\sim w_{\sigma(i)}^\lambda$.
Apply $\alpha$ to the displayed relation and subtract $\lambda$ times
that relation.  This yields
\[
 \sum_j(k_j-\lambda)c_{\tau(j)}=0.
\]
Eliminating $w_s$ exhibits the cone factors as an independent direct sum,
so $k_j\equiv\lambda\pmod{m_j}$ for every $j$.
The local formula of B.12 now proves the assertion on every cone
subgroup and, by conjugation in $\widehat B$, on all its torsion.

\emph{The compact case.}
First suppose that there are at least two cone points.  We use Wilkes's
cohomological calculation with its exact scope:
for a closed orientable hyperbolic orbifold group $B$ with cone
generators $c_j$ of orders $m_j$, an automorphism of $\widehat B$
sending each $c_j$ to a conjugate of $c_j^{k_j}$ acts on
$H^2(\widehat B;\Z/n)$, for every $n$, by one scalar
$\kappa\in\widehat\Z$ satisfying
$\kappa\equiv k_j\pmod{m_j}$.
This is \cite[Proposition 5.7]{bWilkes2017Seifert}, numbered Proposition
6.6 in arXiv:1512.05587v1.  Its proof concerns arbitrary automorphisms
of the indicated form; no lifting to a Seifert group is used in that
calculation.  In particular the scalar is independent of $n$.

For transparency, the part of the calculation that supplies both facts
is as follows.  After taking coinvariants, the relation cells $y_0,y_j$
have boundaries
$\partial y_0=\sum_jx_j$ and $\partial y_j=m_jx_j$.
The automorphism sends $x_j$ to $k_jx_j$, $y_j$ to $k_jy_j$, and
$y_0$ to $\kappa y_0+\sum_j\rho_jy_j$.
Commutation with the boundary gives
$\kappa+m_j\rho_j=k_j$ in $\widehat\Z$.
The dual coboundaries are $y_0^*+m_jy_j^*$, which shows that the
induced action is multiplication by $\kappa$ for every finite
coefficient group.  The completed resolution and the naturality of this
calculation are established in the cited proposition and its use of
Proposition 2.9 there.

The cone permutation of our $\alpha$ is realized by an
orientation-preserving homeomorphism of the underlying surface carrying
the marked cone points to the correspondingly ordered points.  It
induces an orbifold automorphism $\beta$ with the same permutation and
power $1$ on each cone.  Put $\psi=\beta^{-1}\alpha$ and apply the
preceding result to $\psi$.  Let
\[
 \eta_n\in H^2(\widehat B;\Z/n)
\]
be the image of the oriented fundamental cohomology class of the
underlying closed surface under the canonical map from the orbifold
classifying space to its underlying space.  Equivalently, in the
displayed resolution it is $[y_0^*]$, up to the common orientation sign.
The finite-coefficient class is transported to continuous cohomology by
goodness of Fuchsian groups, as in Wilkes's calculation.
Since $\beta$ preserves orientation, $\beta^*\eta_n=\eta_n$.
Thus $\alpha^*\eta_n=\psi^*\beta^*\eta_n=\kappa\eta_n$.
This argument does not require $\beta$ to preserve the chosen subgroup
$P$.

Set $N=[B:\pi]$.  Restriction to the surface subgroup gives
\[
 \operatorname{res}_{P}(\eta_n)=N\omega_n,
 \qquad\omega_n\in H^2(P;\Z/n)
\]
with $\omega_n$ the positive surface class: the underlying map from
the covering surface to the orbifold's underlying surface is an
orientation-preserving branched map of degree $N$.  Naturality of
restriction and A.2's duality yield
\[
 N\kappa\,\omega_n=N\lambda\,\omega_n
       \quad\text{for every }n.
\]
Indeed $\alpha_*[U]=\lambda[U]$ implies
$\langle\alpha^*\omega_n,[U]\rangle
=\langle\omega_n,\alpha_*[U]\rangle=\lambda$.
Taking all finite coefficients gives $N(\kappa-\lambda)=0$ in
$\widehat\Z$, hence $\kappa=\lambda$.  The local congruences are
therefore $k_j\equiv\lambda\pmod{m_j}$, exactly as required.

We remove the single-cone case from the resolution input explicitly.
When $r=1$, the underlying genus $h$ is positive and
$B^{\mathrm{ab}}=\Z^{2h}$, because the one cone generator is zero in
abelianization.  The characteristic subgroup
\[
 B_1=\ker\bigl(B\longrightarrow B^{\mathrm{ab}}/2B^{\mathrm{ab}}\bigr)
\]
has index $2^{2h}$ and contains every conjugate of the cone group.
Its orbifold has $2^{2h}$ cone points of the same order.  The subgroup
$\widehat{B_1}$ and $P\cap\widehat{B_1}$ are $\alpha$-invariant.
The argument for at least two cones applies there, and the finite-cover
calculation at the beginning identifies its surface multiplier with
$\lambda$.  For $r=0$ there is no nonidentity torsion in $\widehat B$,
by D.2.  Thus no exceptional compact signature remains.

Finally, choose positive local cone generators, for which
$\mu(c_j)=1/m_j$ in the chosen orientation line.  The proved congruence
says $\mu(\alpha(c_j))=\lambda/m_j$; its powers have the same
property.  Every positive torsion element is conjugate in $\widehat B$
to one of these powers by D.2.  Together with the reflection case this
proves the stated equation on all torsion of $\widehat E$.
\end{proof}

\noindent\emph{Why the earlier rank-one argument was insufficient.}
Integral relative orbifold cohomology can have a torsion summand:
for $C_3*C_3$ with cusp $\langle ab\rangle$ it has the presentation
$\langle u,v,w\mid3u=w=3v\rangle\cong\Z\oplus\Z/3$.
The proof above uses the cusp relation in the noncompact case and the
action on the full finite-coefficient cohomology in the compact case.
It therefore supplies the missing local-power congruences without
choosing an unproved canonical generator through that torsion summand.

\subsection{Completion of geometric mapping class groups and full loop automorphisms}

\begin{claim}{D.8}{Proved}{All finite subgroups in the geometric completion in the three requested types}
Let $M=\Gamma_{g,\nu}$ include puncture permutations and orientation reversal. For $(g,\nu)=(0,3),(0,4),(1,1)$, completion induces a bijection on conjugacy classes of finite subgroups of $M$ and $\widehat M$. Every finite subgroup of $\widehat M$ is solvable. For every finite $K\leq M$,
\[
 N_{\widehat M}(K)=\overline{N_M(K)},\qquad
 C_{\widehat M}(K)=\overline{C_M(K)}.
\]
These assertions concern $\widehat M$ with its given embedding in $\Out_{\rm lac}(P)$.
\end{claim}
\begin{proof}
For the three-punctured sphere $M=S_3\times C_2$ is finite. For the punctured torus $M=\mathrm{GL}_2(\Z)$; its principal congruence subgroup of level $3$ is a torsion-free free group of finite rank. For the four-punctured sphere the subgroup fixing each puncture and preserving orientation is a free group of rank two and has finite index. Thus all three $M$ are finitely generated virtually free groups.

Bridson--Piwek, Proposition 5.10 \cite{R-BP}, states for every finitely generated virtually free group that finite subgroup conjugacy classes are preserved bijectively by profinite completion. Its proof uses the finite graph of finite groups for $M$: completion gives the profinite fundamental group of the same graph, a finite subgroup fixes a vertex of its profinite Bass--Serre tree, and subgroup conjugacy separability distinguishes the resulting discrete subgroups. This applies without an orientation restriction and proves the first assertion. The normalizer equality is Ribes--Zalesskii, Theorem 2.6 \cite{R-RZ}, for finitely generated subgroups of finitely generated virtually free groups. To deduce the centralizer equality, the homomorphism $N_{\widehat M}(K)\to\Aut(K)$ has finite discrete target. Its kernel is open; the intersection of a dense subgroup with this open kernel is dense in the kernel. This is exactly $C_M(K)$.

For solvability, a finite subgroup of $M_{0,4}$ embeds in $S_4\times C_2$ by the torsion-free pure oriented kernel; one in $\mathrm{GL}_2(\Z)$ embeds in $\mathrm{GL}_2(\F_3)$ by the level-$3$ kernel. The former group is solvable. The latter has normal series with factors inside $C_2$, $A_4$, and the central $C_2$, since $\mathrm{PGL}_2(\F_3)=S_4$; it too is solvable. The finite case $S_3\times C_2$ is immediate. Transfer by the proved conjugacy bijection gives solvability in $\widehat M$.
\end{proof}

\begin{claim}{D.9}{Proved}{The nonsolvable discrete possibilities in genus at most two}
A nonsolvable finite subgroup of the extended mapping class group of a punctured orientable surface of genus at most two is realized on a sphere and has group $A_5$ or $A_5\times C_2$. Its puncture set is invariant under the corresponding icosahedral action. In particular it has at least twelve punctures. This classification does not assert the absence of additional nonsolvable finite subgroups in completions for all larger puncture numbers.
\end{claim}
\begin{proof}
Realize the finite subgroup by Part C. On the sphere its orientation-preserving subgroup is cyclic, dihedral, $A_4$, $S_4$, or $A_5$, by the finite rotation-group classification. The first four are solvable and an extension of one of them by an orientation group of order at most two is solvable. The orthogonal normalizer of the icosahedral rotation group is $A_5\times\{\pm I\}$: a normalizing orthogonal transformation preserves its rotation axes and its icosahedral cellulation, and composition with an icosahedral rotation leaves either $I$ or $-I$. A nonsolvable full group is consequently $A_5$ or $A_5\times C_2$. The smallest nontrivial orbit has twelve vertices; an invariant puncture set has at least twelve points.

On a torus the kernel of the action on $H_1$ consists of a finite translation group, which is abelian. The orientation-preserving linear quotient is cyclic of order $1,2,3,4$, or $6$, and the extended linear quotient is cyclic or dihedral. These groups and their extensions by the translation group are solvable.

In genus two the hyperelliptic involution is central in the full automorphism group of the invariant Riemann surface. Adjoin it when it is absent. The quotient is a sphere with six branch points. A nonsolvable orientation-preserving quotient would contain $A_5$, whose smallest orbit has twelve points, contradicting preservation of the six branch points. The quotient is therefore solvable; so is its central extension by $C_2$ and its extension by orientation reversal. Punctures can only reduce the finite group. This proves the asserted list and its restriction to the discrete setting.
\end{proof}

\begin{claim}{D.10}{Proved}{Finite torsion in the odd-prime normalized pro-$p$ centralizer}
Let $p$ be odd, let $F^{(p)}$ be the free pro-$p$ group on $x,y$, and define $GT_{0,p}$ by the normalized automorphisms
\[
 F_{\lambda,f}(x)=x^\lambda,\qquad
 F_{\lambda,f}(y)=f^{-1}y^\lambda f,
 \quad \lambda\in\Z_p^\times,\quad f\in(F^{(p)})',
\]
satisfying the two-cycle and three-cycle equations. Every finite subgroup of $GT_{0,p}$ is cyclic of order dividing $p-1$. There is a subgroup of order $p-1$, and every finite subgroup is conjugate in $GT_{0,p}$ to a subgroup of it. For $p\geq5$ this supplies finite loop-preserving outer automorphism groups which are not conjugate to geometric groups of the three-punctured sphere.
\end{claim}
\begin{proof}
The normalized description is faithful: a conjugator between two normalized automorphisms first belongs to $C(x)=\langle x\rangle$; the normalization $f\in F'$ and the $y$ equation remove its exponent. The same observation identifies these automorphisms with the pure outer centralizer of the three puncture permutations. The equations are closed conditions, composition is the usual substitution law, and their inverses satisfy the same equations. Hence this is a profinite group. The Galois action on $\mathbf P^1_{\overline\Q}-\{0,1,\infty\}$ gives elements of this group with multiplier the $p$-adic cyclotomic character. This character surjects onto $\Z_p^\times$: the finite cyclotomic extensions have Galois groups $(\Z/p^n)^\times$, and their inverse limit is $\Z_p^\times$.

The kernel of $\Aut(F^{(p)})\to\Aut(F^{(p)}/\Phi(F^{(p)}))$ is pro-$p$; the exact statement and its application are also given in Herfort--Zalesskii, proof of Theorem 6.3 \cite{R-HZ}. For the normalized group the action on this Frattini quotient is the scalar $\lambda\bmod p$. Thus the kernel of
\[
 GT_{0,p}\longrightarrow\F_p^\times
\]
is pro-$p$, and the map is surjective. The profinite Schur--Zassenhaus theorem supplies a complement $C_{p-1}$ and conjugacy of complements in the inverse image of any subgroup of $\F_p^\times$. One can see the passage from the finite theorem directly: apply the finite theorem to every quotient by an open normal subgroup of the pro-$p$ kernel, then take the inverse limit of the nonempty finite sets of compatible complements and conjugators. The finite intersection property follows by passing to a common finite quotient.

We check the necessary absence of torsion in the kernel of the multiplier. A free pro-$p$ group has its injective Magnus map into
$\Z_p\langle\!\langle X,Y\rangle\!\rangle^\times$, with $x\mapsto1+X$, $y\mapsto1+Y$. The map is injective because reduction modulo $p$ is the usual faithful free pro-$p$ Magnus map; it extends continuously from the free group for the $(p,X,Y)$-adic topology on the power-series algebra. The ordinary homogeneous degree is used only in the following leading-term calculation, not as the topology for that extension. An automorphism acting identically on $F^{(p)}_{\rm ab}=\Z_p^2$ acts identically on the degree-one quotient of this power-series algebra. For a nonidentity such automorphism choose the least degree $d\geq2$ of a nonzero correction to $X$ or $Y$. The degree-$d$ correction to its $n$th iterate is $n$ times that correction: substitution of the higher terms can change only degrees above $d$. Each homogeneous coefficient module is a torsion-free $\Z_p$-module, so no positive iterate is the identity. Consequently the multiplier is injective on a finite subgroup.

The torsion subgroup of $\Z_p^\times$ is cyclic of order $p-1$. Therefore every finite subgroup maps injectively to $\F_p^\times$. It is a complement in the inverse image of its image, so the preceding conjugacy theorem proves the full assertion. Finally the geometric pure group of type $(0,3)$ consists of $1$ and complex conjugation, with multipliers $1$ and $-1$. A generator with multiplier a primitive $(p-1)$st root cannot be conjugated into it; the multiplier is unchanged by conjugation even in the full loop-preserving outer group. The first example has order four at $p=5$.
\end{proof}

\begin{claim}{D.11}{Proved}{Executed finite checks and their exact scope}
The program \texttt{code/root/finite\_checks.py} verifies the two-cycle and three-cycle equations in every Heisenberg group $H(\F_p)$ for $p=5,7,11$ and every $\lambda\in\F_p^\times$, using
\[
 f(a,b)=[a,b]^{(\lambda^2-1)/24},\qquad
 m=(\lambda-1)/2.
\]
In particular it verifies an order-four normalized automorphism at $(p,\lambda\bmod p)=(5,2)$ and an order-three one at $(7,2)$. The multipliers of the corresponding infinite finite-order subgroups are Teichm\"uller lifts of these residues; the integer $2\in\Z_5^\times$ itself has infinite order. D.10, not the finite checks alone, proves the existence of compatible infinite pro-$p$ lifts.
\end{claim}
\begin{proof}
The program defines all $p^3$ elements as triples and uses
$(a,b,c)(a',b',c')=(a+a',b+b',c+c'+ab')$.
It enumerates inverses and powers and evaluates each side of both equations in that group. Its output records every multiplier and both truth values. Since all commutators are central, the associated automorphism is
$(a,b,c)\mapsto(\lambda a,\lambda b,\lambda^2c)$; the program checks multiplication preservation and computes its order. Division by $24$ is only used at the three displayed primes. No assertion of this formula at $2$ or $3$ is made. At $p=3$, D.10 gives only order-two prime-to-$p$ torsion. At $p=2$, the Magnus argument still proves that a finite subgroup injects into $\{\pm1\}$, but the complement argument does not settle conjugacy of involutions in $GT_{0,2}$.
\end{proof}

\begin{claim}{D.12}{Open}{Torsion in the full profinite $GT_0$ and the full low-type groups}
D.29 refutes the finite geometric-recognition assertion for $(g,\nu)=(0,4)$ by an actual involution with a complete preserved peripheral system. We do not settle whether every nonidentity torsion element of the full profinite $GT_0$ is conjugate to $(-1,1)$, or the full finite-subgroup recognition questions for $(g,\nu)=(0,3),(1,1)$. The geometric-completion assertion in all three low types is proved in D.8, and the odd-prime analogue is decided in D.10. D.32 proves that every finite subgroup of the full pure three-cusp group, including $GT_0$, has order at most two; it gives the faithful map $G\to S_3\times C_2$ for the full three-cusp group. D.34, D.36, D.41 and D.42 identify the global orientation character of every such finite group: every involution has its global $2$-adic sign, and every order-three element has multiplier $1$. D.35 identifies the exact remaining third-cusp conjugator equation, proves full generation from it, and identifies its conjugator sets in the pure group and in $GT_0$. D.38 gives a complete characteristic $A_5^{19}$ quotient test, including proved exclusions of finite candidates from full lifts. D.44 solves the coboundary and the first symmetry equation simultaneously in the full group; the third-cusp equation and hence existence of the conjugator for each actual full action remain unresolved. D.33 proves that full homology detects every finite subgroup in type $(1,1)$ and that its order divides $24$. D.43 proves that every involution in this full punctured-torus group has its global orientation sign. D.45 identifies the exact characteristic binary-icosahedral refinement and its full finite solution counts; the same cusp-triangle survivors pass. D.46 proves that every cyclic order-four outer action on the punctured-torus group has an actual order-four lift, using the three full section classes of its positive square. D.47 checks the exact conjugator coordinates in the binary refinement: precisely 10,125,000 of the 324,000,000 coordinate pairs pass, with the count unchanged by finite invertibility and both symmetry relations. D.48 proves that a full punctured-torus order-three action has exactly two section classes and a torsion-free cusp normalizer whose cube exponent is a 3-adic unit. D.49 computes the standard pure involution centralizer and the exact two-element fibres of its conjugacy orbit map. D.50 classifies the finite two-subgroups of the full punctured-torus group as exactly $1,C_2,C_4,C_2\times C_2,D_8$, exhibits all five, and proves the global multiplier character on every four-group or dihedral group of order eight. D.51 identifies the second torsion-free normal index-three kernel of every order-three torus extension, with its exact type-$(1,3)$ intersection cover and three labelled peripheral normalizers. D.52 gives the twelve possible abstract finite-subgroup types, realizes nine of them, and proves the reflection multiplier for every subgroup of type $S_3,D_{12},D_{24}$. D.53 proves the unique full labelled order-four section class and its exact normalizer. D.54 excludes every nonzero prime-to-three invariant character for a cubic torus action by an explicit peripheral norm comparison in a cyclic cover. D.55 supplies the corresponding square comparison, proves that every positive involution acts as $-I$ at every prime, and proves faithfulness of every primary homology map on every finite subgroup. The finite torus abstract types are consequently exactly the nine realized types $1,C_2,C_3,C_4,C_2\times C_2,C_6,S_3,D_8,D_{12}$; the three extra types in the earlier allowed list are excluded. D.56 computes the prime-section permutations of every cyclic-six action and proves that its full extension has no order-six subgroup or finite-order representative of the outer generator. D.57 proves actual labelled lifts for every four-group and order-eight dihedral action, with the exact three-or-one and one full section counts. D.58 determines the exact negative-involution section counts and their full fixed-subgroup mod-two cohomology, and proves that the two internal cusp-reflection classes fuse under full conjugacy. D.59 identifies the negative cusp section as the class with zero full fixed-subgroup mod-two first cohomology, proves that its entire fixed group has zero ambient abelian image, and proves that every actual cusp-reflection transporter generates the full fixed homology line. Its odd cyclic covers have exactly the displayed cusp cycles and full section counts, and the other fixed group, when present, surjects onto that homology line. D.60 identifies the quotient by the closed normal closure of the cusp reflection with the exact full homology coinvariants, and the quotient by all torsion with $\widehat\Z$. It constructs a section in $F$, and in the two-class case a section inside the other fixed group, giving an actual embedded $\widehat\Z\times C_2$. D.61 proves that the peripheral commutator belongs to the closed derived subgroup of no proper closed subgroup of the full rank-two free group. It consequently gives exact negative-conjugator parameter spaces with no additional generation condition, retaining both cusp-reflection indices and all full unit multipliers. Existence of these parameters, pure-involution orbit exhaustion, full fixed-group recognition, the remaining scalar cubic torus multipliers, and geometric conjugacy remain unresolved; D.70 below settles the quartic multiplier.

D.62 gives the exact full negative-conjugator transporter cosets. D.63--D.65 compute the full standard negative centralizers and their homological characters. D.66 computes the exact scalar-cover norm preimage and primary abelianized-extension splittings. D.67 gives the standard Klein-group and dihedral normalizers and all finite overgroups of the specified two-section negative involution. D.68 identifies the unique invariant double cover of every actual quartic action, its integral homology maps, and the full invariant dihedral quotient. D.69 determines the exact rotation image of the positive-square fixed subgroup in that finite quotient, its two labelled section classes, its fixed-kernel first cohomology and the full integral positive-square action on the nontrivial deck summand. D.70 uses the actual closed filling and its pro-two quotient to exclude every negative odd quartic coordinate by a degree-three restriction contradiction. Hence every order-four full torus element has global multiplier one, and every finite two-subgroup has its multiplier in the two global signs. D.71 constructs the actual cubic Frobenius cover, proves the full integral regular module on its nontrivial deck summand, and determines the image and first mod-three cohomology of the relevant section fixed subgroup. D.72 gives an explicit surjective quasi-isomorphism from the actual normalized pants cochain algebra to the three-section Boolean graded algebra; its acyclic kernel proves the stated strong Massey vanishing and exact finite unitriangular representation image. The arbitrary negative conjugator equations, scalar cubic multiplier and full low-type recognition remain unresolved.
D.73 proves the exact compatible $q$-adic cubic Frobenius quotient and its full fixed-group splitting, the integral module and section constraints at every $q$-power level, the limiting fixed-kernel character vanishing, and an actual two-torsion-generator obstruction. D.74 proves that all three full pants section fixed subgroups are trivial. Its actual finite-coefficient relative duality, full Brown maps and coset-module witness give self-centralizing sections and real projectivity. The full coboundary $h$ in D.35 is consequently unique, including D.44's first symmetry. The exact third-cusp equation with this $h$, arbitrary negative torus conjugacy, and the scalar cubic realization or exclusion remain unresolved.
D.75 proves full triviality of every negative torus cusp-reflection fixed subgroup, uniqueness of the actual cusp-fusion transporter, and both singleton transporter formulas for every full unit. In the one-section-class case its extension is real projective. Its exact remaining equation is the conjugacy (D.75.5); the other fixed group in the two-section-class case retains its full homology surjection without a procyclic recognition conclusion.
D.76 identifies the actual pants section augmentation module with $\mathbb F_2[[E]]^2$, using all finite coefficient modules and compatible finite-quotient isomorphisms. Its specified incidence map has cokernel equal to the augmentation kernel of $\mathbb F_2[[E/H]]$, where $H$ is the closed subgroup generated by the three specified section representatives. Thus its invertibility records exactly $H=E$. Four explicit involutions generate $E$, with the duplicated class and the resulting three-involution subgroup both written out. Generation by the specified triple and the original third-cusp equation remain unresolved.
D.77 constructs a complete three-class involution triple in the standard full pants extension whose two kernel generators give automorphisms of every maximal pro-$p$ quotient, while an actual $A_5$ quotient sends the triple to $D_{10}$ and its kernel to $C_5$. Its specified incidence cokernel maps onto a five-dimensional augmentation module. The same quotient proves failure of conjugacy to every unit power of the original second cusp. That retained cusp condition prevents this example from deciding the distinguished triple of D.76 or the equation (D.74.40).
D.78 proves that every open subgroup of the other negative torus fixed group has mod-two first cohomology of dimension one. Its specified dyadic procyclic subgroup is a Sylow subgroup, and the full fixed group has the actual split coordinates $B=K\rtimes\widehat{\mathbb Z}$ with a pro-odd kernel $K$. The finite-coefficient inflation and section-restriction maps retain $Q^K\hookrightarrow Q$. The original finite-coset invariant map detects every nonidentity element of $K$; its surjectivity and full procyclic recognition remain unresolved.
\end{claim}
Here is the exact obstruction reached, rather than an extrapolation from computation. The homomorphism from a free profinite group to the product of its maximal pro-$p$ quotients does not detect all its finite quotients. Our program constructs $A_5$ from a $5$-cycle and a $3$-cycle, enumerates its sixty elements, and verifies that its commutator subgroup has sixty elements. Thus $A_5$ is a nontrivial two-generated finite quotient invisible to every nilpotent quotient. The pro-$p$ complements constructed in D.10 do not give a compatible solution of the defining equations in such quotients.

We also checked the $f=1$ attack: the three-cycle relation becomes the power-product relation of Lochak--Schneps, Lemma 9 \cite{R-LS}; their finite modular quotients force the global alternatives $\lambda=1$ or $-1$. A choice of independent signs in the factors of $\widehat\Z$ is not a solution of this relation. The finite nilpotent search therefore supplies no full-profinite counterexample. D.32 excludes every element of odd prime order in the full pure three-cusp group and every element of finite order greater than two. The smallest unresolved full-$GT_0$ case is therefore an additional involution up to conjugacy; D.34 proves that its multiplier is globally $-1$, supplies its three section classes and a commutator coboundary, and leaves the conjugator inside the specified loop group to be determined. D.23 supplies the exact order-bounded compatibility criterion over all isolated finite shadows. D.24 executes nonnilpotent refinement tests and proves that specific order-three shadows cannot lift with order three; D.25 computes the entire dihedral inverse limit and the exact residual multiplier-one norm equation. D.28 proves the additional prime-order section and fixed-subgroup constraint, including its exact maximal pro-$p$ comparison. The finite-cover argument of D.32 now excludes all odd-order full lifts, beyond those bounded computations; the involution conjugacy problem retains its exact isolated-refinement formulation. D.29 uses an arbitrary Galois pants isomorphism and its inverse on two different pieces; this produces a finite involution in type $(0,4)$ without requiring that pants isomorphism itself have finite order. For genus at most two and arbitrary puncture number, the solvable comparison is supplied by the checked argument in F.7. D.31 now proves that every nonsolvable finite subgroup in these completions is abstractly $A_5$ or $A_5\times C_2$ on a punctured sphere, and gives exactly the geometric positive puncture-orbit sizes. Its proof also excludes every nonsolvable finite subgroup in the genus-one and genus-two completions. What remains is geometric conjugacy of the actual $A_5$ embeddings and the corresponding centralizer comparison of F.7.

The later tests also close two specific attempted constructions of a lower-type counterexample. D.39 constructs an actual rank-two Haran involution but proves, using its unique section class and the trivial centralizer of that section, that it preserves neither hyperbolic loop marking under any full profinite change of basis. D.37 gives the exact four-cusp descent condition for a negative transposition. D.40 excludes the required additional deck symmetry for the actual D.29 example by the unequal homomorphism counts $10200$ and $4920$ in $\operatorname{PSL}_2(\F_{11})$; the argument also excludes deck maps moving the original separator. These are proved obstructions for the specified constructions. They do not provide the conjugator still required in D.35.

\subsection{Descent along a finite covering}

Put $A(P)=\Aut_{\rm lac}(P)$ and let a discretization mean the image of a loop-preserving isomorphism $\widehat{\pi_0}\to P$ restricted to $\pi_0$. Thus it is an actual dense subgroup with its peripheral classes. Define $\operatorname{Discr}(P)$ using actual subgroups, and define $\operatorname{Discr}^{\rm out}(P)$ by quotienting this set by conjugation by $P$. This convention is necessary for an outer action: an outer automorphism acts canonically on the second set. Write $J=\Aut_{\rm lac}(\pi_0)$, $B=\Out_{\rm lac}(P)$ and $M=\Out_{\rm lac}(\pi_0)$ for their embedded images.

\begin{claim}{D.13}{Proved}{Restriction and normal descent with the peripheral step included}
For an open subgroup $V\leq P$, restriction of the corresponding finite covering defines an injective map
\[
 \operatorname{Discr}(P)\longrightarrow\operatorname{Discr}(V),
 \qquad\Delta\longmapsto\Delta\cap V.
\]
For normal $V$, its image consists exactly of the lattices $\Lambda\subset V$ whose outer deck action $P/V\to\Out_{\rm lac}(V)$ is induced by orientation-preserving loop automorphisms of $\Lambda$. The cyclic peripheral subgroups of the descended lattice have closures equal to the specified peripheral subgroups of $P$.
\end{claim}
\begin{proof}
We prove the self-normalizer input, including the class-preserving step. Paris \cite[Theorems 1.7 and 3.2]{R-Paris} proves conjugacy $p$-separability for closed orientable surface groups and free groups. Thus two elements of $\Lambda$ conjugate in $\widehat\Lambda$ are already conjugate in $\Lambda$.

For a free group $F_r=\langle x_1,\ldots,x_r\rangle$, let $\varphi$ preserve every conjugacy class. After an inner adjustment we have $\varphi(x_1)=x_1$. Write $\varphi(x_i)=w_ix_iw_i^{-1}$. Choose a reduced representative $w$ for the double coset $\langle x_1\rangle w_i\langle x_i\rangle$, so that $w_i=x_1^awx_i^b$ and $w$ begins with neither $x_1$ nor $x_1^{-1}$ and ends with neither $x_i$ nor $x_i^{-1}$. The word $x_1wx_iw^{-1}$ is cyclically reduced with length $2+2|w|$. Since $\varphi(x_1x_i)$ is conjugate to $x_1x_i$, we obtain $w=1$. Hence
$\varphi(x_i)=x_1^{n_i}x_ix_1^{-n_i}$. For distinct $i,j\geq2$, conjugating $\varphi(x_ix_j)$ by $x_1^{-n_i}$ gives
$x_i x_1^{n_j-n_i}x_jx_1^{n_i-n_j}$; cyclic length is two exactly when $n_i=n_j$. Thus all $n_i$ agree, and $\varphi$ is inner. The rank-two case has only one exponent and gives the same conclusion.

For a closed orientable surface group of genus $g\geq2$, Dehn--Nielsen--Baer realizes the outer class of a class-preserving automorphism by a mapping class $f$. It fixes every oriented closed-curve class and acts identically on $H_1$, so preserves orientation. Choose a finite filling Alexander system of simple curves. The Alexander method (\cite[\S2.3]{R-FarbMargalit}) makes the subgroup fixing all those curve classes finite: its possible permutations of intersection points and edges form a finite group; the kernel fixes the resulting graph and is isotopic to the identity on its complementary discs. Therefore $f$ has finite order. Realize the cyclic group generated by $f$ using Part C. For an action of order $d$, rational cochains of the quotient are the invariant cochains upstairs, by an equivariant triangulation and averaging over the finite stabilizers. The quotient's first Betti number is consequently $2g$, because the action on $H_1$ is trivial. Its underlying genus is $g$. Riemann--Hurwitz now gives
\[
 2-2g=d\left(2-2g-\sum_j(1-1/m_j)\right).
\]
For $d>1$ the right side is strictly smaller than $2-2g<0$, a contradiction. Thus $f=1$ and the automorphism is inner.

Now let $h\in N_{\widehat\Lambda}(\Lambda)$. Its conjugation automorphism is class-preserving by the first paragraph, hence is conjugation by an element of $\Lambda$ by the two arguments just given. Dividing $h$ by that element centralizes the dense lattice and therefore all of $\widehat\Lambda$. Part A gives trivial centre, so $h\in\Lambda$. This proves $N_{\widehat\Lambda}(\Lambda)=\Lambda$ with the exact comparisons needed in LM3 B.8.

For normal $V$, a descended lattice $\Delta$ normalizes $\Lambda=\Delta\cap V$. Conversely, an element of $N_P(\Lambda)$ has a coset in the finite group $P/V$. The intersection $N_P(\Lambda)\cap V$ is $\Lambda$ by the just-proved input. Given that the deck action is induced on $\Lambda$, form the pullback discrete extension $E_\Lambda$ of $P/V$ by $\Lambda$ inside $\Aut_{\rm lac}(\Lambda)$. We construct its map into $P$ explicitly. An element of the pullback is a pair $(g,\varphi)$, with $g\in P/V$, $\varphi\in\Aut_{\rm lac}(\Lambda)$, and prescribed outer class. Extend $\varphi$ continuously to $\widehat\varphi\in\Aut(V)$ and choose a lift $p_g\in P$ of $g$. Equality of outer classes gives a $v\in V$ with $\widehat\varphi=\Inn(vp_g)|_V$. This $v$ is unique, since $C_P(V)=1$ by Part A. Send $(g,\varphi)$ to $vp_g$. Changing $p_g$ changes $v$ inversely and leaves the product unchanged; multiplication is respected because both products induce the same conjugation on $V$ and have the same quotient label. This defines an injective homomorphism: an element in its kernel has $g=1$ and induces the identity on the dense subgroup $\Lambda$. Its image is exactly $N_P(\Lambda)$, by restricting a normalizing conjugation to $\Lambda$ and applying the same uniqueness. The discrete group is torsion-free, as a subgroup of $P$. The finite-extension realization theorem of Part C gives a surface lattice $\Delta=E_\Lambda$; its orientation is positive because the deck action on $V$ has multiplier $+1$. Its completion is $P$ by D.1.

We supply the peripheral argument left incomplete in LM3 B.10. Let $C\leq\Lambda$ be a primitive peripheral subgroup, and let $Q=C_P(\overline C)$ be the maximal procyclic peripheral subgroup in $P$ containing $\overline C$. Put $D=C_\Delta(C)$. The quotient surface supplied by Part C makes $D$ its primitive cusp subgroup. We have $\overline D\leq Q$ and
\[
 \overline D\cap V=\overline C=Q\cap V.
\]
The image of $D$ in $P/V$ is the stabilizer of the $\Lambda$-conjugacy class of $C$. To see this, adjust a representative of a stabilizing coset by an element of $\Lambda$ to normalize $C$; orientation preservation changes normalization into centralization. Conversely a centralizing element fixes the class. The identical argument in the profinite group identifies the image of $Q$ with the stabilizer of the $V$-conjugacy class of $\overline C$. The map from discrete cusp classes to their profinite classes is bijective and equivariant by peripheral noncommensurability and the finite covering construction of Part A. These two stabilizers are therefore equal. The two closed subgroups $\overline D\leq Q$ have the same intersection with the open subgroup $V$ and the same image in $P/V$, so they are equal. This proves the promised loop identification.

In particular a descended lattice is forced to equal $N_P(\Lambda)$, proving injectivity for normal $V$. For arbitrary $V$, take $W=\bigcap_{p\in P}pVp^{-1}$, an open normal subgroup, and use the actual intersection map $\Lambda\mapsto\Lambda\cap W$. Equality after restriction to $V$ implies equality after restriction to $W$, hence equality of the original lattices by the normal case. This proves injectivity without assuming $V$ normal.
\end{proof}

\begin{claim}{D.14}{Proved}{The image for a nonnormal finite covering}
Let $V\leq P$ be open, $W=\operatorname{core}_P(V)$ and $\Lambda\in\operatorname{Discr}(V)$. Set $\Lambda_W=\Lambda\cap W$. The image condition for $\Lambda$ is that the full deck action of $P/W$ on $W$ is induced by orientation-preserving loop automorphisms of $\Lambda_W$. This condition is independent of a choice of coset representatives. The unique extension is $N_P(\Lambda_W)$, and its intersection with $V$ is $\Lambda$.
\end{claim}
\begin{proof}
Necessity follows by restriction of an actual lattice in $P$. For sufficiency apply the normal descent statement D.13 to $W\triangleleft P$, obtaining $\Delta=N_P(\Lambda_W)$. Both $\Delta\cap V$ and $\Lambda$ are discretizations of $V$ with intersection $\Lambda_W$ with the normal open subgroup $W$ of $V$. The injectivity already proved inside $V$ gives $\Delta\cap V=\Lambda$. The deck action is an outer action, so changing a representative multiplies its automorphism by an inner automorphism of $W$; its membership in the embedded geometric outer subgroup is unaffected. This proves the asserted independence and the exact nonnormal image description.
\end{proof}

\begin{claim}{D.15}{Proved}{The explicit nonextending discretization in scan 354}
Let $U=\mathbf P^1-\{0,1,\infty\}$ and
$U'=\mathbf P^1-\{0,\infty,1,-1\}$, with the double cover
$z\mapsto z^2$. Put $P=\widehat{\pi_1(U)}$, let $V\triangleleft P$ be its index-two subgroup, and let $\Lambda=\pi_1(U')\subset V$ be the geometric discretization. There is an explicit continuous loop automorphism $\gamma$ of $V$ for which $\gamma(\Lambda)$ is a discretization of $V$ that does not extend to $P$, although it has the same prediscretization as $\Lambda$.
\end{claim}
\begin{proof}
The deck involution $s:z\mapsto-z$ fixes $0,\infty$ and exchanges $1,-1$. Let $a$ be a simple curve separating $\{0,1\}$ from $\{\infty,-1\}$ and let $b=s(a)$; write $T_a,T_b$ for positive Dehn twists. They give a free basis of the pure oriented mapping class group $M^{!+}_{0,4}\simeq F_2$: in the ordered moduli space $\mathcal M_{0,4}\simeq\mathbf P^1-\{0,1,\infty\}$ they are monodromies around two distinct boundary points. One can verify the completed embedding directly from A.12: filling the fourth puncture identifies this pure group with the point-pushing group $\pi_{0,3}$, and the peripheral-normalization map sends its completion into $\Aut(\widehat\pi_{0,3})$ as $\Inn(\widehat\pi_{0,3})$. The latter is isomorphic to $\widehat\pi_{0,3}$ by triviality of its centre. The composite is injective, so the completion embeds already in $\Out_{\rm lac}(V)$. Thus the pure group's closure is its free profinite completion, with abelianization $\widehat\Z^2$ on the indicated twists.

Choose the fully specified profinite integer
\[
 \alpha_2=0\in\Z_2,\qquad
 \alpha_\ell=1\in\Z_\ell\quad(\ell\ {\rm odd}),
 \qquad \alpha\in\widehat\Z.
\]
It is not an integer: its $2$-adic coordinate would force that integer to be zero, contrary to its $3$-adic coordinate. Let $u=T_a^\alpha$ in the closure of $M^{!+}_{0,4}$. Since $s$ preserves orientation,
$sT_as^{-1}=T_b$, and the continuous power maps give
\[
 u^{-1}su=s\,T_b^{-\alpha}T_a^\alpha .
\]
The pure factor on the right has abelianization $(\alpha,-\alpha)$. Every element of the discrete pure group has integral abelianization. Moreover
$M_{0,4}\cap\overline{M^{!+}_{0,4}}=M^{!+}_{0,4}$: the puncture permutation and orientation sign are detected by finite quotients. Consequently $u^{-1}su$ does not belong to the discrete extended mapping class group.

Choose any continuous automorphism $\gamma\in\Aut_{\rm lac}(V)$ representing $u$. Its ambiguity is conjugation by an element of $V$, which does not change membership in the outer geometric subgroup. The geometric outer subgroup associated to $\gamma(\Lambda)$ is $uM_{0,4}u^{-1}$. The calculation says that the deck involution is absent from this subgroup. The necessary direction of D.13 therefore proves that $\gamma(\Lambda)$ does not extend to $P$. Since $u$ lies in the closure of the pure geometric group, its coset modulo the closed geometric stabilizer is the identity, proving equality of the prediscretizations.

The code in \texttt{code/root/finite\_checks.py} verifies compatible Chinese-remainder residues of $\alpha$ and the puncture partitions interchanged by $s$. Nonintegrality and the infinite completed-group conclusion follow from the coordinate argument above. This spells out and verifies the proposed element in LM3's review, \S4.2.
\end{proof}
\noindent\emph{Manuscript:} \S28, folios 178--179, scans 352--354, especially \emph{Un tel $\gamma$ existe toujours}. We have proved the indicated index-two example; the word \emph{toujours} is not used to assert the construction for every finite covering.


\begin{claim}{D.16}{Proved}{A precise fixed-set formula for every actual action}
Let $A$ be a group, $J\leq A$, and $G\leq A$ a finite subgroup. Under the left action of $G$ on $A/J$,
\[
 (A/J)^G=\{aJ:a^{-1}Ga\leq J\}.
\]
It is the disjoint union of the $C_A(G)$-orbits
\[
 \coprod_{[a]} C_A(G)\big/\bigl(aC_J(a^{-1}Ga)a^{-1}\bigr),
\]
where $[a]$ ranges over $J$-conjugacy classes of the labelled embeddings $g\mapsto a^{-1}ga$ occurring in $J$. This formula applies to actual discretizations with $(A,J)=(A(P),\Aut_{\rm lac}(\pi_0))$, to outer discretizations with $(A,J)=(B,M)$, and to prediscretizations with $(A,J)=(B,\overline M)$.
\end{claim}
\begin{proof}
The fixed-coset condition is $gaJ=aJ$ for every $g$, exactly the displayed inclusion. Two fixed cosets $aJ,bJ$ are in the same centralizer orbit exactly when there are $c\in C_A(G)$ and $j\in J$ with $b=caj$. Substitution gives $b^{-1}gb=j^{-1}a^{-1}gaj$ for every labelled $g$. Conversely this equality makes $b j^{-1}a^{-1}$ centralize $G$. The stabilizer of $aJ$ is $C_A(G)\cap aJa^{-1}=aC_J(a^{-1}Ga)a^{-1}$, proving the orbit formula.

The homogeneous-space identifications use the extension of each discrete loop automorphism to completion and its injectivity. The stabilizer of an actual lattice in $A(P)$ is exactly $J$. Its inner-conjugacy quotient is $B/M$. Since the closure of $\Inn(\pi_0)$ is $\Inn(P)$, the closed based stabilizer contains all inner automorphisms and its quotient is $\overline M$; hence the prediscretization set is $B/\overline M$. These identifications also show explicitly why a bare outer element is not assigned an action on the set of actual subgroups without choosing a lift.
\end{proof}

\subsection{Torsion-free extensions and recognition}

\begin{claim}{D.17}{Proved}{Torsion-free extensions with an invariant discretization}
Let $1\to P\to\mathcal E\to G\to1$ be a profinite extension with $G$ finite, $P=\widehat\pi_0$ hyperbolic, and the outer action preserving the loop structure and a discretization. A torsion-free $\mathcal E$ is the profinite completion of a torsion-free discrete surface group, with the descended peripheral system. The quotient surface may be nonorientable.
\end{claim}
\begin{proof}
First the outer action is faithful. Its kernel is the image of $C_{\mathcal E}(P)$: an element acting inner can be divided by its unique inner element of $P$. The centralizer meets $P$ trivially, so embeds in the finite quotient $G$ and is finite. Torsion-freeness makes it trivial. Construct the discrete extension $E$ of the stabilized lattice with this faithful outer action, using centre-freeness. By D.1, its completion is the given $\mathcal E$ with the identity on kernel and quotient. The discrete extension embeds in its completion because it is the finite extension realized by Part C and hence residually finite. Torsion-freeness of $\mathcal E$ therefore implies torsion-freeness of $E$. In a realization, a nontrivial point stabilizer gives a nontrivial finite-order element of the lift group. Thus $G$ acts freely on the surface associated to $\pi_0$, and its quotient is a surface. Its fundamental group is $E$ by lifting paths. The primitive cusp group in $E$ is the centralizer of the covering cusp group; D.13 proves equality of their closures with the specified procyclic peripherals. Orientation-reversing free deck transformations yield the nonorientable quotient and its orientation character. This is the exact scope of the descent argument.
\end{proof}

\begin{claim}{D.18}{Proved}{The unmarked pro-$p$ extension statement for punctured surfaces}
A finitely generated torsion-free pro-$p$ group $\mathcal E$ containing an open free pro-$p$ group $P$ of rank $r$ is free pro-$p$. Its rank is
\[
 1+\frac{r-1}{[\mathcal E:P]}.
\]
In particular the displayed quantity is an integer. This proves the underlying-group assertion for punctured pro-$p$ surfaces, without identifying the peripheral system. More generally, every torsion-free profinite group with an open normal pro-$p$ subgroup is itself pro-$p$. Thus D.19--D.20 apply to arbitrary finite-quotient extensions of their pro-$p$ kernels; the quotient is forced to be a $p$-group.
\end{claim}
\begin{proof}
Herfort--Zalesskii, Theorem 1.1 \cite{R-HZ}, writes such a finitely generated virtually free pro-$p$ group as the fundamental pro-$p$ group of a finite graph of finite $p$-groups, with vertex groups embedded in $\mathcal E$. Torsion-freeness makes all vertex and edge groups trivial. The fundamental pro-$p$ group is therefore the free pro-$p$ group on the cycles of the graph. The pro-$p$ Schreier rank formula gives $r-1=[\mathcal E:P](\operatorname{rank}(\mathcal E)-1)$, proving the formula. The theorem retains no designated peripheral subgroups in its graph-of-groups conclusion; their identification is a further assertion, not supplied by this calculation.

We prove the assertion about a general finite quotient. Write $G=\mathcal E/P$ and fix a prime $q\ne p$. For every open normal $U\triangleleft\mathcal E$ contained in $P$, the kernel of $\mathcal E/U\to G$ is a finite $p$-group. A Sylow $q$-subgroup of $\mathcal E/U$ therefore has order $|G|_q$ and maps isomorphically to a Sylow $q$-subgroup of $G$. Fix one such subgroup $Q\leq G$. The finite nonempty sets of Sylow $q$-subgroups above $Q$ admit compatible choices under quotient maps: an image is again a Sylow subgroup, and any specified one lifts by taking a Sylow subgroup of its inverse image. Compactness gives an inverse-limit subgroup of $\mathcal E$ isomorphic to $Q$, since all these transition maps on the chosen subgroups are isomorphisms. Torsion-freeness forces $Q=1$. Hence every prime divisor of $|G|$ is $p$, and every finite quotient of $\mathcal E$ is a $p$-group. This is exactly the assertion that $\mathcal E$ is pro-$p$.
\end{proof}

\begin{claim}{D.19}{Proved}{Recognition of the peripheral pro-$p$ extension}
Let $p$ be any prime, let $P$ be the pro-$p$ completion of an orientable hyperbolic surface group of type $(g,\nu)$ with $\nu>0$, and let
\[
 1\longrightarrow P\longrightarrow\mathcal E\longrightarrow G\longrightarrow1
\]
be an exact sequence of pro-$p$ groups with $G$ finite. Assume that $\mathcal E$ is torsion-free and its outer action preserves the $P$-conjugacy classes of the peripheral subgroups. Choose one $L_j$ from each $\mathcal E$-orbit of these classes. Then
\[
 \mathcal N_j=N_{\mathcal E}(L_j)=C_{\mathcal E}(L_j)\simeq\Z_p,
\]
and the pair $(\mathcal E;\{\mathcal N_j\})$ is continuously isomorphic, preserving the indexed peripheral conjugacy classes, to the pro-$p$ completion of a compact surface with its boundary circles. Its interior is orientable for odd $p$. At $p=2$ the interior may be nonorientable; its orientation character is the induced action on the orientation line of $P$.
\end{claim}
\begin{proof}
Each $\mathcal N_j$ is closed, and
$\mathcal N_j\cap P=N_P(L_j)=L_j$ by Part A. Thus $\mathcal N_j$ is a torsion-free pro-$p$ group with an open subgroup $\Z_p$. It is finitely generated, by lifting generators of its finite quotient and adjoining a generator of $L_j$. Apply D.18 to $\mathcal N_j$: the free pro-$p$ rank formula gives rank one. Consequently $\mathcal N_j\simeq\Z_p$, and it centralizes $L_j$.

We identify the restricted peripheral family, including its multiplicities. The restriction of $(\mathcal E;\{\mathcal N_j\})$ to $P$ contains the subgroups
\[
 P\cap e\mathcal N_je^{-1}=eL_je^{-1},
 \qquad e\in P\backslash\mathcal E/\mathcal N_j .
\]
The map from this double-coset set to the orbit of the $P$-conjugacy class of $L_j$ is surjective by definition. It is injective: an equality
$e'L_je'^{-1}=a eL_je^{-1}a^{-1}$ with $a\in P$ is equivalent to
$e'^{-1}ae\in N_{\mathcal E}(L_j)=\mathcal N_j$, exactly equality of the double cosets. Thus the restricted family consists of all original peripheral classes, each once.

Wilkes, \cite[Proposition 5.11]{R-WilkesRelative}, proves finite-index ascent and descent of Poincar\'e duality pairs at $p$ for a $p$-torsion-free profinite ambient group, with compatible dualizing modules. The pair $(P;\{L_i\})$ is orientable of dimension two: this follows from its standard free pro-$p$ surface presentation and the converse direction of \cite[Theorem 3.3]{R-Wilkes}. We have checked the exact restricted family and torsion hypothesis. Hence $(\mathcal E;\{\mathcal N_j\})$ is a Poincar\'e duality pair of dimension two.

Its dualizing module has underlying group $\Q_p/\Z_p$, with action through a continuous character $\chi:\mathcal E\to\Z_p^\times$ whose restriction to $P$ is trivial. Naturality of duality under conjugation identifies its character with the action on the surface orientation line, using the inverse convention on the dual cohomology line. The image of $\chi$ is a finite $p$-subgroup because it factors through $G$: it is trivial at odd $p$ and is trivial or $\{\pm1\}$ at $p=2$. Thus inversion of the character does not affect the answer. Its restriction to each $\mathcal N_j$ is trivial, by the boundary-duality restriction in \cite[Proposition 5.9]{R-WilkesRelative} and the trivial duality character of $\Z_p$.

D.18 makes $\mathcal E$ free pro-$p$. Write the peripheral count as $b+1$ and the rank as $n+b$. The condition $n\geq0$ in the classification follows from the relative cohomology sequence: the cokernel of
\[
 H^1(\mathcal E;\F_p)\longrightarrow
       \bigoplus_{j=0}^{b}H^1(\mathcal N_j;\F_p)
\]
has dimension one, since $H^2(\mathcal E;\F_p)=0$ and the relative top group has dimension one. Hence $b\leq\operatorname{rank}(\mathcal E)$.
The classification in \cite[Theorem 3.3]{R-Wilkes} gives a free basis
$x_1,\ldots,x_n,s_1,\ldots,s_b$, with each $s_i$ generating a conjugate of $\mathcal N_i$, and a generator $s_0$ of $\mathcal N_0$ in one of these forms:
\[
\begin{array}{ll}
 \chi=1:&
 s_0=[x_1,x_2]_{-}\cdots[x_{n-1},x_n]_{-}\,s_1\cdots s_b,
 \\[2pt]
 \chi\ne1,\ n\ {\rm odd}:&
 s_0=x_1^2[x_2,x_3]_{-}\cdots[x_{n-1},x_n]_{-}\,s_1\cdots s_b,
 \\[2pt]
 \chi\ne1,\ n\ {\rm even}:&
 s_0=x_1^2[x_1,x_2]_{-}[x_3,x_4]_{-}\cdots
 [x_{n-1},x_n]_{-}\,s_1\cdots s_b .
\end{array}
\]
Here $[x,y]_{-}=x^{-1}y^{-1}xy$ matches that source's convention. The first line has $n$ even. The other two have $p=2$; their parameter $2^f$ is zero, since $\chi$ has no infinite principal-unit part.

These are surface words, with explicit changes of basis. Inverting each handle generator changes the commutator convention in the first line. The second line gives one crosscap and orientable handles. For the third, put
\[
 a=x_1x_2^{-1},\qquad b'=x_2,\qquad
 x_1=ab',\quad x_2=b'.
\]
Then $x_1^2[x_1,x_2]_{-}=a^2(b')^2$, which gives two crosscaps and orientable handles. Replacing $s_0$ by its inverse gives the usual boundary relation and preserves its subgroup. Sending the discrete surface generators to these free pro-$p$ generators induces the asserted isomorphism of completed pairs, with exactly the indexed peripheral conjugacies already specified. The free rank formula gives
$\chi(\mathcal E)=\chi(P)/[\mathcal E:P]<0$, so its interior is hyperbolic. This proves the peripheral statement at every prime, including $2$ and $3$.
\end{proof}

\begin{claim}{D.20}{Proved}{Recognition of the closed pro-$p$ extension}
Let $P$ be the pro-$p$ completion of a closed orientable surface group of genus $g\geq2$. A torsion-free pro-$p$ group $\mathcal E$ containing $P$ as an open normal subgroup is the pro-$p$ completion of a closed surface group. The surface is orientable at odd $p$. At $p=2$ a nonorientable surface is allowed, and its orientation character is the action on the orientation line of $P$.
\end{claim}
\begin{proof}
We give the absolute finite-index duality calculation. The orientable surface presentation makes $P$ a Poincar\'e duality pro-$p$ group of dimension two, with
\[
 H^i(P,\Z_p[[P]])=
 \begin{cases}\Z_p,&i=2,\\0,&i\ne2,\end{cases}
\]
and trivial right $P$-action on the top group. This is the trivial-character case of the Demushkin classification in \cite[Theorem 1.9]{R-Wilkes}.

Serre's theorem \cite[\S1, Theorem]{R-SerreCD} equates the $p$-cohomological dimensions of an open subgroup and a profinite group without elements of order $p$. Therefore $\operatorname{cd}_p(\mathcal E)=2$. Wilkes \cite[Proposition 1.26]{R-WilkesRelative} gives finite-index ascent and descent of type $p$-$FP_\infty$ for profinite groups. The surface resolution gives this property for $P$, and hence for $\mathcal E$.

At finite index, induction and coinduction of compact modules agree, and continuous Shapiro gives
\[
\begin{split}
 H^i(\mathcal E,\Z_p[[\mathcal E]])
 &=
 H^i\!\left(\mathcal E,
        \operatorname{coind}_{P}^{\mathcal E}\Z_p[[P]]\right)\\
 &\simeq H^i(P,\Z_p[[P]]).
\end{split}
\]
The first coefficient isomorphism and compact Shapiro are \cite[Propositions 1.16 and 1.21, and the remark after 1.21]{R-WilkesRelative}. We make their module compatibility explicit. Put $\Lambda_H=\Z_p[[H]]$ and let $\operatorname{pr}_P:\Lambda_{\mathcal E}\to\Lambda_P$ retain the $P$-coset component. The coefficient isomorphism is
\[
 \Lambda_{\mathcal E}\longrightarrow\operatorname{Coind}_P^{\mathcal E}\Lambda_P,
 \qquad \lambda\longmapsto \bigl(x\mapsto\operatorname{pr}_P(x\lambda)\bigr).
\]
Here coinduction consists of left-$P$-linear functions and its left $\mathcal E$-action is $(e f)(x)=f(xe)$. For representatives $t$ of the finite set $P\backslash\mathcal E$, the inverse is $f\mapsto\sum_t t^{-1}f(t)$. These continuous maps are left-$\mathcal E$-linear and right-$P$-linear. Evaluation at $1$ in Shapiro is natural for the right-$P$ coefficient endomorphisms. Consequently the top right duality character restricts to that of $P$. The cohomology criterion for Poincar\'e duality now makes $\mathcal E$ a duality pro-$p$ group of dimension two with dualizing module $\Q_p/\Z_p$. Its orientation character $\chi:\mathcal E\to\Z_p^\times$ has finite $p$-group image, trivial at odd $p$ and contained in $\{\pm1\}$ at $p=2$.

Its one-dimensional degree-two cohomology and nondegenerate degree-one cup product make it an infinite Demushkin group. The classification \cite[Theorem 1.9]{R-Wilkes} gives, for these characters, exactly the three words of D.19 with the boundary letters omitted. The same invertible free-basis change converts the even-rank nonorientable word to the surface form. Hence an actual closed surface presentation induces a continuous isomorphism of its pro-$p$ completion with $\mathcal E$. Finite-index multiplicativity for the Euler characteristic of the finite two-dimensional resolution gives
\[
 \chi(\mathcal E)=\frac{2-2g}{[\mathcal E:P]}<0.
\]
The resulting surface is hyperbolic, and conjugation-naturality of the dualizing module gives the stated orientation.
\end{proof}

\begin{claim}{D.21}{Refuted}{A full profinite normalizer shortcut}
The assertion that every torsion-free profinite group with an open subgroup $\widehat\Z$ is itself $\widehat\Z$ is false. The pro-$p$ version used in D.19 is proved in D.18.
\end{claim}
\begin{proof}
Define
\[
 N=(\Z_3\rtimes\Z_2)\times\prod_{\ell\ne2,3}\Z_\ell,
 \qquad (a,b)(a',b')=(a+(-1)^{\,b\bmod2}a',\,b+b')
\]
on the first two factors, and use addition on the remaining factors.
This is a profinite group with the open normal subgroup
\[
 L=(\Z_3\times2\Z_2)\times\prod_{\ell\ne2,3}\Z_\ell
       \simeq\widehat\Z
\]
of index two. A finite-order element projects trivially to the torsion-free quotient
$\Z_2\times\prod_{\ell\ne2,3}\Z_\ell$, and then lies in the torsion-free kernel $\Z_3$; it is therefore the identity. Yet conjugation by $(0,1)$ sends $(1,0)$ to $(-1,0)$, so $N$ is nonabelian. The code in \texttt{code/root/finite\_checks.py} checks its quotient
$C_3\rtimes C_4$, the index-two cyclic subgroup of order six, and this noncommutation. The proof in the inverse limit establishes torsion-freeness; no finite quotient is claimed to be torsion-free.

The multiplier on $L$ is $-1$ in its $3$-adic coordinate and $+1$ in every other coordinate. Thus passing separately to pro-$p$ groups loses precisely this interaction between two different primes. D.29 uses this $N$ as an explicit retract in its full peripheral counterexample. Its cusp normalizers are all procyclic, so that example does not depend on making $N$ itself a cusp normalizer.
\end{proof}

\begin{claim}{D.22}{Refuted}{Full profinite recognition without an invariant discretization}
The assertion that every torsion-free finite extension of a full
profinite hyperbolic surface group with preserved loop structure
is a completed surface pair is false.  D.29 constructs a
counterexample with kernel of type $(0,4)$ and quotient $C_2$,
including the exact peripheral comparisons and a finite quotient
excluding every free profinite presentation. D.30 constructs a
counterexample with closed genus-two kernel and quotient $C_2$;
its degree-two cohomology at two odd primes excludes every
surface completion. Recognition with an
invariant discretization holds by D.17.  The separate pro-$p$
recognition theorems, with their precise group and peripheral
hypotheses, are proved in D.18--D.20.
\end{claim}
\begin{proof}
The basis-derived four-cusp presentation, faithful outer involution,
continuous twisted extension and torsion-free retraction are given
in D.29.  Its order-eighteen quotient needs three generators, whereas
an index-two extension of a free rank-three group would have rank two
if it were free.  The exact finite-cover comparison excludes closed
surface completions.  Thus the counterexample satisfies the full
peripheral hypothesis and fails the claimed conclusion.

For the closed case, D.30 gives the extension by its continuous
twisted multiplication and by an equivalent profinite presentation.
Its torsion-free maximal pro-$2$ quotient detects the parity of
every element, proving that every torsion element lies in its
torsion-free surface kernel. The two restriction isomorphisms give
$H^2(\mathcal E;\mathbf F_3)=\mathbf F_3$ and
$H^2(\mathcal E;\mathbf F_5)=0$. The cohomology of every orientable,
nonorientable, or punctured surface completion is computed in
D.30 and excludes this pair. This settles the closed hyperbolic
assertion in its smallest genus as well.

This also identifies the failed step in the earlier recognition
attempt.  Separate pro-$p$ comparisons do not produce a single
full profinite surface presentation.  The mixed-prime multiplier
in D.29 is an actual finite-order loop automorphism, and all its
cusp normalizers are procyclic.  The preserved global surface
relation therefore does not restore the proposed conclusion.
The first possible punctured nonabelian rank is three: D.27 excludes
proper torsion-free finite extensions of $\widehat F_2$.
The explicit free-basis comparison in D.29 settles the formerly
undecided peripheral marking on D.26's rank-three extension.

We checked Collas--Philip--Yamaguchi \cite{R-CPY}, Definition 2.1,
Theorems 1.3 and 2.4, against the supplied source.  Their profinite
F-groups already have the stated orbifold presentation.
Theorem 2.4 constructs a torsion-free open normal subgroup, and
Theorem 1.3 bounds the derived length of a finite quotient in its
nonperfect case.  The maps go from a prescribed presentation to
an open subgroup; they do not assert the reverse reconstruction
from an arbitrary finite extension.  D.29 refutes that reverse
assertion under the present full profinite hypotheses.  The
positive conclusions of D.17--D.20 retain their proved maps and
hypotheses throughout.
\end{proof}

